\documentclass{article}%
\usepackage[T1]{fontenc}%
\usepackage[utf8]{inputenc}%
\usepackage{lmodern}%
\usepackage{textcomp}%
\usepackage{lastpage}%
\usepackage{geometry}%
\geometry{margin=1in,headheight=14pt,headsep=25pt}%
\usepackage{hyperref}%
\usepackage{enumitem}%
\usepackage{fancyhdr}%
\usepackage{xcolor}%
\usepackage{url}%
\usepackage{breakurl}%
\usepackage{amsmath}%
\usepackage{amssymb}%
\usepackage{mathtools}%
\usepackage{amsthm}%
\usepackage{thmtools}%
%

            \hypersetup{
                colorlinks=true,
                linkcolor=blue,
                filecolor=magenta,
                urlcolor=blue,
                breaklinks=true,
                bookmarks=true,
                pdfborder={0 0 0}
            }
            \urlstyle{same}
            \Urlmuskip=0mu plus 1mu  % URL breaking in nicer spots
            
            % Math configuration
            \everymath{\displaystyle}
            \setlength{\jot}{10pt}  % Increase spacing between equations
            
            \pagestyle{fancy}
            \fancyhf{}
            \rhead{Generated on \today}
            \cfoot{\thepage}
            
            % Better Unicode support
            \usepackage[utf8]{inputenc}
            \usepackage{newunicodechar}
            \usepackage[breaklinks=true]{hyperref}
            \setlength{\itemsep}{1em} % Adds spacing between items
            
            % Configure itemize settings globally
            \setlist[itemize]{nosep}
            \setlist[itemize]{leftmargin=*}
            
            % Remove section numbering
            \renewcommand{\thesection}{}
            \renewcommand{\thesubsection}{}
        %
\lhead{Lecture: 2024{-}09{-}24{-}SimplexMatrix}%
\title{2024{-}09{-}24{-}SimplexMatrix}%
\author{Generated by Scribe.AI}%
\date{\today}%
%
\begin{document}%
\normalsize%
\maketitle%
\section{Slide 1}%
\label{sec:Slide1}%
\subsection{Content}%

\underline{\text{Simplex in Matrix Notation}} \\
\begin{array}{|c|}
\hline
\text{max } f(x) = c^T x \\
\text{s.t. } Ax \le b \\
x \ge 0 \\
\hline
\end{array}
\quad x \in \mathbb{R}^n, \quad c \in \mathbb{R}^n \\
A \in \mathbb{R}^{m \times n}, \quad b \in \mathbb{R}^m \\
\begin{array}{|c|}
\hline
\text{max } f(x) = c^T x \\
\text{s.t. } w = b - Ax \\
x, w \ge 0 \\
\hline
\end{array}
\quad w \in \mathbb{R}^m
%
\subsection{Description}%
This slide presents two equivalent formulations of a linear program in matrix notation. The title "Simplex in Matrix Notation" is underlined. The first formulation is a standard maximization problem: maximize $f(x) = c^Tx$ subject to $Ax \le b$ and $x \ge 0$, where $x \in \mathbb{R}^n$, $c \in \mathbb{R}^n$, $A \in \mathbb{R}^{m \times n}$, and $b \in \mathbb{R}^m$. The second formulation introduces slack variables $w \in \mathbb{R}^m$, rewriting the constraints as $w = b - Ax$ and $x, w \ge 0$.  The equivalence between the two formulations is implied by a curved arrow connecting them.  The context of the course (Linear Programming and Optimization Techniques) suggests that this slide introduces the matrix notation used to represent linear programs, which is essential for understanding the simplex method and other optimization algorithms. The slide likely serves as an introduction to the matrix representation of linear programs before delving into the simplex algorithm.%
\subsection{Figures}%
\begin{figure}[h]%
\centering%
\begin{minipage}{\textwidth}%
\centering%
\begin{minipage}{0.15\textwidth}%
\centering%
\includegraphics[width=0.8\textwidth]{/Users/ashoksaravanan/Coding/ScribeLec/server/output/MA 421/lectures/2024-09-24-SimplexMatrix/figures/1.2.png}%
\caption{First linear program in matrix notation.}%
\end{minipage}%
\hfill%
\begin{minipage}{0.15\textwidth}%
\centering%
\includegraphics[width=0.8\textwidth]{/Users/ashoksaravanan/Coding/ScribeLec/server/output/MA 421/lectures/2024-09-24-SimplexMatrix/figures/1.1.png}%
\caption{Second linear program in matrix notation.}%
\end{minipage}%
\end{minipage}%
\end{figure}

%
\newpage%
\section{Slide 2}%
\label{sec:Slide2}%
\subsection{Content}%

\underline{\text{Simplex in Matrix Notation}} \\
\begin{array}{|c|}
\hline
\text{max } f(x) = c^T x \\
\text{s.t. } Ax \le b \\
x \ge 0 \\
\hline
\end{array}
\quad x \in \mathbb{R}^n, \quad c \in \mathbb{R}^n \\
A \in \mathbb{R}^{m \times n}, \quad b \in \mathbb{R}^m \\
\begin{array}{|c|}
\hline
\text{max } f(x) = c^T x \\
\text{s.t. } w = b - Ax \\
x, w \ge 0 \\
\hline
\end{array}
\quad w \in \mathbb{R}^m \\
\begin{array}{|c|}
\hline
\text{max } f(x) = j^T + c^T x \\
\text{s.t. } \overline{w} = \overline{b} - \overline{A}x \\
x, \overline{w} \ge 0 \\
\hline
\end{array}
%
\subsection{Description}%
This slide shows three equivalent formulations of a linear program in matrix notation. The title "Simplex in Matrix Notation" is underlined. The first formulation is a standard maximization problem: maximize $f(x) = c^Tx$ subject to $Ax \le b$ and $x \ge 0$, where $x \in \mathbb{R}^n$, $c \in \mathbb{R}^n$, $A \in \mathbb{R}^{m \times n}$, and $b \in \mathbb{R}^m$. The second formulation introduces slack variables $w \in \mathbb{R}^m$, rewriting the constraints as $w = b - Ax$ and $x, w \ge 0$. The third formulation is similar to the second, but uses different notation, with $\overline{w}$, $\overline{b}$, and $\overline{A}$ representing modified slack variables, right-hand side vector, and constraint matrix, respectively, and includes a constant term $j^T$ in the objective function.  Curved arrows indicate the iterative process of the simplex method, transforming the problem from one formulation to the next. The context of the course (Linear Programming and Optimization Techniques) suggests that this slide illustrates how linear programs are represented in matrix form, which is crucial for understanding the simplex algorithm.  The slide likely serves as a transition to explaining the simplex method's iterative process.%
\subsection{Figures}%
\begin{figure}[h]%
\centering%
\begin{minipage}{\textwidth}%
\centering%
\begin{minipage}{0.15\textwidth}%
\centering%
\includegraphics[width=0.8\textwidth]{/Users/ashoksaravanan/Coding/ScribeLec/server/output/MA 421/lectures/2024-09-24-SimplexMatrix/figures/2.3.png}%
\caption{First linear program in matrix notation.}%
\end{minipage}%
\hfill%
\begin{minipage}{0.15\textwidth}%
\centering%
\includegraphics[width=0.8\textwidth]{/Users/ashoksaravanan/Coding/ScribeLec/server/output/MA 421/lectures/2024-09-24-SimplexMatrix/figures/2.2.png}%
\caption{Second linear program in matrix notation.}%
\end{minipage}%
\hfill%
\begin{minipage}{0.15\textwidth}%
\centering%
\includegraphics[width=0.8\textwidth]{/Users/ashoksaravanan/Coding/ScribeLec/server/output/MA 421/lectures/2024-09-24-SimplexMatrix/figures/2.1.png}%
\caption{Third linear program in matrix notation.}%
\end{minipage}%
\end{minipage}%
\end{figure}

%
\newpage%
\section{Slide 3}%
\label{sec:Slide3}%
\subsection{Content}%

\underline{\text{Simplex in Matrix Notation}} \\
\begin{array}{|c|}
\hline
\text{max } f(x) = c^T x \\
\text{s.t. } Ax + w = b \\
x, w \ge 0 \\
\hline
\end{array}
\quad \begin{array}{c} x \in \mathbb{R}^n \\ c \in \mathbb{R}^n \\ A \in \mathbb{R}^{m \times n} \\ b \in \mathbb{R}^m \\ w \in \mathbb{R}^m \end{array}
\begin{array}{c} m \times (n+m) \\ \begin{bmatrix} A & I \end{bmatrix} \begin{bmatrix} x \\ w \end{bmatrix} = b \end{array} \\
\xrightarrow{\text{Simplex iterations}} \\
\begin{array}{|c|}
\hline
\text{max } f(x) = \overline{j} + \overline{c}^T x \\
\text{s.t. } \overline{A}x + \overline{w} = \overline{b} \\
x, \overline{w} \ge 0 \\
\hline
\end{array}
\quad \begin{array}{c} m \times (n+m) \\ \begin{bmatrix} \overline{A} & I \end{bmatrix} \begin{bmatrix} x \\ \overline{w} \end{bmatrix} = \overline{b} \end{array}
%
\subsection{Description}%
This slide illustrates the simplex method's iterative process using matrix notation within the context of linear programming. The title "Simplex in Matrix Notation" is underlined.  The slide presents two equivalent formulations of a linear program. The first is a standard maximization problem: maximize $f(x) = c^Tx$ subject to $Ax + w = b$ and $x, w \ge 0$, where $x \in \mathbb{R}^n$ are the decision variables, $c \in \mathbb{R}^n$ is the cost vector, $A \in \mathbb{R}^{m \times n}$ is the constraint matrix, $b \in \mathbb{R}^m$ is the right-hand side vector, and $w \in \mathbb{R}^m$ are the slack variables. This is then represented in matrix form as $\begin{bmatrix} A & I \end{bmatrix} \begin{bmatrix} x \\ w \end{bmatrix} = b$. The second formulation represents a later iteration of the simplex method, where the objective function and constraints have been modified. The objective function is now $\overline{j} + \overline{c}^T x$, and the constraints are $\overline{A}x + \overline{w} = \overline{b}$, with the same non-negativity constraints. This is also represented in matrix form as $\begin{bmatrix} \overline{A} & I \end{bmatrix} \begin{bmatrix} x \\ \overline{w} \end{bmatrix} = \overline{b}$. The arrow indicates the iterative nature of the simplex algorithm, transforming the initial problem into a modified one during each iteration. The dimensions of the matrices are explicitly stated. The slide emphasizes the use of matrix notation to represent and solve linear programs efficiently.%
\subsection{Figures}%
\begin{figure}[h]%
\centering%
\begin{minipage}{\textwidth}%
\centering%
\begin{minipage}{0.15\textwidth}%
\centering%
\includegraphics[width=0.8\textwidth]{/Users/ashoksaravanan/Coding/ScribeLec/server/output/MA 421/lectures/2024-09-24-SimplexMatrix/figures/3.4.png}%
\caption{First linear program in matrix notation.}%
\end{minipage}%
\hfill%
\begin{minipage}{0.15\textwidth}%
\centering%
\includegraphics[width=0.8\textwidth]{/Users/ashoksaravanan/Coding/ScribeLec/server/output/MA 421/lectures/2024-09-24-SimplexMatrix/figures/3.1.png}%
\caption{Second linear program in matrix notation.}%
\end{minipage}%
\hfill%
\begin{minipage}{0.15\textwidth}%
\centering%

\includegraphics[width=0.8\textwidth]{/Users/ashoksaravanan/Coding/ScribeLec/server/output/MA 421/lectures/2024-09-24-SimplexMatrix/figures/3.3.png}%
\caption{Matrix representation of the first linear program.}%
\end{minipage}%
\end{minipage}%
\end{figure}

%
\newpage%
\end{document}